\documentstyle[%


righttag%


,11pt%


,amssymb%


,russian%               remove in Eng.


%,izvru%                 Set izven% in Eng.


,izvru-ks%             for brief communications


]{amsart}





%       Math definitions





\newcommand{\vraisup}{\operatornamewithlimits{vrai\,sup}}


\newcommand{\vsup}[2]{\vraisup\limits_{#2}{#1}}


\newcommand{\eqdef}{\overset{\mathrm{def}}{=}{}}


\newcommand{\seta}{\longrightarrow}


\newcommand{\tts}[1]{}





\theoremstyle{plain}


\theorembodyfont{\it}


\newtheorem{thmnonum}{\hskip\parindent Теорема}


\renewcommand{\thethmnonum}{}





%\hoffset=-15mm%                  For Eng.


\setcounter{page}{56}


\god{2000}


\nomer{4~(455)} %                        Remove in Eng.


%\nomer{Vol.\,44, No.\,4}%               Set in Eng.


%\str{\pageref{begin}--\pageref{end}}%  Set in Eng.


\udk{517.929}





\author{\it М.Ж.\,Алвеш}


% NB:   ^^ remove in Eng.


\title{Об одной нелинейной краевой задаче\\


с несуммируемой особенностью}





\thanks{Работа выполнена при финансовой поддержке


Capacity Building Project (Credit 2436, Eduardo Mondlane University


--- Mozambique).}





\date{}


%\pagestyle{myheadings}%         Set in Eng.


\pagestyle{plain} %              Remove in Eng.


\begin{document}


%\thispagestyle{plain}%          Set in Eng.


\maketitle


\label{begin}





\section{Пространство ${B}_p^v$}





Обозначим через


${C}$ пространство непрерывных функций


$x:[0,1]\seta{\bold R}^1$ с нормой


$\|x\|_{C}\eqdef\max\limits_{0\le t\le 1}|x(t)|$, через


${L}_p$ ($1\le p<\infty$) --- пространство классов эквивалентных


суммируемых в $p$-й степени функций $x:[0,1]\seta{\bold R}^1$ с


нормой $\|x\|_{L_p}\eqdef\Big (\int\limits_0^1


|x(t)|^p\,dt \Big)^{1/p}$, а через ${L}_\infty$ ---


пространство классов эквивалентных измеримых и ограниченных в


существенном функций $x:[0,1]\seta{\bold R}^1$ с нормой


$\|x\|_{L_\infty}\eqdef\vsup{|x(t)|}{0\le t\le 1}$.


Пусть $v :[0,1]\seta [0,+\infty]$ --- неубывающая положительная на


$(0,1]$ функция.





\begin{defn}[{\series{m}\shape{n}\selectfont\cite{1}}]


${B}^v_p$


($1\le p < \infty$) --- пространство классов эквивалентных


измеримых функций $x:[0,1]\seta {\bold R}^1$ таких, что


для каждой функции $x\in { B}^v_p$ существует константа


$M_x \ge 0$,


для которой всюду на $[0,1)$


выполнено неравенство


$\int\limits_0^t |x(s)|^p\,ds \le M^p_x\,v(t)$.


\end{defn}





В качестве нормы


$\|x\|_{B^v_p}$ элемента $x\in { B}^v_p$


можно взять величину


$\|x\|_{B^v_p}\eqdef \inf M_x$. Норму


$\|x\|_{B^v_p}$ элемента $x\in{ B}^v_p$


можно также определять как


$\|x\|_{B^v_p}\eqdef


\sup\limits_{0<t<1}{\Big(\frac{1}{v(t)}\int\limits_0^t


|x(s)|^p\,ds\Big)^{1/p}}$. При\linebreak $p=\infty$ через


${B}^v_\infty$


определим пространство классов эквивалентных


измеримых функций $x:[0,1]\seta {\bold R}^1$ таких, что


для каждой функции $x\in { B}^v_p$ существует константа


$M_x \ge 0$,


для которой всюду на $(0,1)$


выполнено неравенство


$\vsup{|x(s)|}{0\le s\le t} \le M_x\,v(t)$. Норму


$\|x\|_{B^v_\infty}$


элемента $x\in { B}^v_\infty$


определяем выражением


$\|x\|_{B^v_\infty}\eqdef \inf M_x$. Норму


$\|x\|_{B^v_\infty}$ элемента $x\in { B}^v_\infty$


можно также определять как


$\|x\|_{B^v_\infty}\eqdef


\sup\limits_{0<t<1}


{(\frac{1}{v(t)}\vsup{|x(s)|}{0\le s\le t})}$.


Пространства $B^{v}_p$ банаховы.





\begin{lem}


Пусть пространства


$B^{v_1}_p$ и $B^{v_2}_p$


порождены функциями $v_1$ и $v_2$


соответственно, и пусть имеет место неравенство


$M^p_{12} \eqdef \sup\limits_{0<t<1}


{\frac{v_1(t)}{v_2(t)}} <\infty$.





Тогда пространство $B^{v_1}_p$


непрерывно вложено в пространство $B^{v_2}_p$,


причем для любого $x\in B^{v_1}_p$


имеем


$\|x\|_{B^{v_2}_p} \le M_{12}\|x\|_{B^{v_1}_p}$.


\end{lem}





Введем обозначение $t^*=\sqrt[\uproot{2}2p]{t}$, если $1 \le p < \infty$, и


$t^*=1$, если $p = \infty$.


Далее введем {\it весовое} пространство ${ B}^{t,t^*}_p$ таких


классов эквивалентных


измеримых функций $x: [0,1]\seta {\bold R}^1$, что $x\in{ 


B}^{t,t^*}_p$, если $xt^*\in{ B}^{t}_p$. Положим


$\|x\|_{B^{t,t^*}_p}\eqdef \|xt^*\|_{B^t_p}$.


Пространство ${ B}^{t,t^*}_p$ банахово.





\section{Линейная сингулярная задача в


пространстве ${D}^{t,t^*}_p$}





Через ${ D}^{t,t^*}_p$ обозначим пространство


функций


$x:[0,1]\seta {\bold R}^1$


таких, что


{\it


$1)$~$x$ непрерывна на $[0,1]$, причем $x(0) = 0;$


$2)$~$\Dot x$ абсолютно непрерывна на $[0,1];$


$3)$~$\Ddot x$ принадлежит пространству $B^{t,t^*}_p$.


}





Это пространство алгебраически изоморфно пространству


${B}_p^{t,t^*}\times {\bold R}^1$. Действительно,


изоморфизмы


${\cal J}:{ B}^{t,t^*}_p\times{\bold R}^1 \seta


{ D}^{t,t^*}_p$


и


${\cal J}^{-1}: { D}^{t,t^*}_p


\seta { B}^{t,t^*}_p\times{\bold R}^1$


можно определить равенствами


${\cal J}=\{\Lambda, Y\}$,


${\cal J}^{-1}=[\delta, r]$,


где


$$


(\Lambda z)\eqdef\int_0^t (t-s)z(s)\,ds,\quad


(Y\beta)\eqdef\beta t,\quad


\delta x \eqdef\Ddot x,\quad


r x\eqdef \Dot x(0).


$$





Если в пространстве ${ D}^{t,t^*}_p$


определить норму равенством


$\|x\|_{D^{t,t^*}_p}\eqdef


\|\Ddot x\|_{B^{t,t^*}_p}+


|\Dot x(0)|$,


то ${ D}^{t,t^*}_p$ --- банахово пространство.





Рассмотрим линейную краевую задачу


\begin{equation}


({\cal L}_0 x)(t)\eqdef \Ddot x(t)+\frac{\Dot x(t)}{t}


- \frac{x(t)}{t^2}=f(t),


\quad t\in [0,1],\quad


x(0)=0,\quad x(1) + \nu \Dot x(1) = \alpha,


\label{1}


\end{equation}


где $\nu \ge 0$.





Значение оператора ${\cal L}_0$ на функции, которая


отлична от нуля в точке $t=0$ (напр., $x(t)\equiv 1$),


не принадлежит пространству суммируемых функций. Поэтому уравнение


${\cal L}_0 x=f$


естественно рассматривать в пространстве, где таких функций нет.


Задачу (\ref{1}) будем исследовать


в пространстве ${ D}^{t,t^*}_p


\simeq { B}^{t,t^*}_p\times {\bold R}^1$.


Для


простоты вычислений ограничимся случаем $\nu\equiv 0$.





{\it Главная часть} (\cite{2}, с.\,103) $Q_0$ оператора


${\cal L}_0: { D}^{t,t^*}_p\seta { B}^{t,t^*}_p$


имеет вид


$$


(Q_0 z)(t)\eqdef ({\cal L}_0\Lambda z)(t)=z(t)+


({\cal K} z)(t),\quad t\in [0,1],


$$


где $({\cal K}z)(t)


\eqdef\frac{1}{t^2}\int\limits_0^t s\,z(s)\,ds$ --- оператор


со ``средним весом'' \cite{3}.


При $1<p\le\infty$ оператор


${\cal K}: { L}_p\seta{ L}_p$


ограничен, причем


$\|{\cal K}\|_{L_p\to L_p}\le


\frac{p'}{2}$ \cite{3}, где $p'$ --- сопряженный показатель.


Более того, при $1 \le p \le\infty$


оператор


${\cal K} : { B}^{t,t^*}_p\seta { B}^{t,t^*}_p$


ограничен, причем


$$


\|{\cal K}\|_{B^{t,t^*}_p\to B^{t,t^*}_p}


\le (1/2)^{1/p'}< 1,\ \; \text{если} \ \; 1 < p \le \infty,


\ \; \text{и} \ \;


\|{\cal K}\|_{B^{t,t^*}_1\to B^{t,t^*}_1}


\le 1.


$$





Оператор ${Q}_0: { B}^{t,t^*}_p\seta { B}^{t,t^*}_p$


при $1 < p \le \infty$ имеет ограниченный обратный


$$


({Q}^{-1}_0z)(t)=z(t) - \frac{1}{t^3}\int_0^t


s^2\,z(s)\,ds.


$$





Таким образом, {\it главная краевая задача} (\cite{2}, с.\,104)


${\cal L}_0 x=f$, $\Dot x(0)=\alpha$


однозначно разрешима.


Фундаментальная система решений однородного уравнения


${{\cal L}}_0 x = 0$


в пространстве ${ D}^{t,t^*}_p$


одномерна и состоит из функции


$x_1(t) = t$.


Итак, главная краевая задача интегрируется в


конечном виде и ее общее решение имеет представление


$$


x(t)=({\Lambda}{Q}_0^{-1}f)(t) + \Dot x(0) t =


\int_0^t (t-s)\{f(s)


- \frac{1}{s^3}\int_0^s \tau^2\,f(\tau)\,d\tau


\}\,ds + \Dot x(0) t.


$$





После перемены порядка интегрирования это представление принимает


канонический вид


$$


x(t)=\int_0^t {C}_0(t,s)f(s)\,ds + \Dot x(0) t=


({\cal C}_0 f)(t) + \Dot x(0) t,


$$


где функция Коши определяется равенством


${C}_0(t,s)= \frac{t^2 - s^2}{2t}$,


$0\le s\le t\le 1$.


Оператор Коши ${\cal C}_0:{ B}_p^{t,t^*}\seta { C}$


вполне непрерывен при $p=\infty$.


В соответствии с ([2], с. 105) для любого функционала $\ell$,


определенного на ${ D}_p^{t,t^*}$


и такого, что


$\ell [x_1]\ne 0$, краевая задача


${{\cal L}}_0x =z$, $\ell x = \alpha$ однозначно разрешима.


Решение задачи (\ref{1}) при $\nu = 0$ в пространстве


${ D}^{t,t^*}_p$


имеет представление


$x={\cal W}f + \alpha x_1$,


где


$({\cal W}f)(t)\eqdef\int\limits_0^1 W(t,s)f(s)\,ds$,


$$


{W}(t,s)\eqdef


\begin{cases}


s^2(t^2-1)/(2t),& \text{если } 0\le s\le t \le 1;\\


t(s^2-1)/2,& \text{если } 0\le t < s \le 1.


\end{cases}


$$


При $ p=\infty$


оператор ${\cal W}: B_p^{t,t^*}\seta { {C}}$


вполне непрерывен.





Рассмотрим задачу


\begin{equation}


({\cal L} x)(t)\eqdef ({\cal L}_0 x)(t) - (Tx)(t)=f(t),\quad


t\in [0,1], \quad


x(0)=0,\quad x(1) = \alpha,


\label{2}


\end{equation}


где


$T : { C}\seta { B}^{t,1}_{\infty}$ --- линейный антитонный


оператор, $f\in { B}_{\infty}^{t,1}$, $\alpha\in {\bold R}^1$.





Пусть ${\cal A}\eqdef {\cal W}T :{ C}\seta { C}$.


Имеет место





\begin{lem}


Следующие


утверждения эквивалентны$:$


\begin{itemize}


\item[1)] существует такой элемент


$y\in D^{t,1}_{\infty}$, что


$y(t) > 0$, $\phi (t)\eqdef ({\cal L}_0 y)(t)-(Ty)(t)\le 0$, $t\in


(0,1)$,


причем $y(1) -\int\limits_0^1\phi (s)\, ds>0;$


\item[2)] спектральный радиус $\rho ({\cal A})$ оператора


${\cal A} : C\seta C$


меньше единицы$;$


\item[3)] краевая задача $(\ref{2})$ имеет единственное решение


$x\in D^{t,1}_{\infty}$


для каждой пары $f\in B^{t,1}_{\infty}$, $\alpha\in


\bold R^1$, причем ее оператор Грина ${\cal G}$ антитонен$;$


\item[4)] существует положительное на $[0,1]$ решение однородного


уравнения ${\cal L}_0 x - Tx=0$.


\end{itemize}


\end{lem}





\section{Нелинейная задача с несуммируемой особенностью}





Пусть $v,z \in { D}^{t,1}_{\infty}$, где $v(t)\le z(t)$


при всех $t \in [0,1]$. Обозначим


$$


[v,z]_{D^{t,1}_{\infty}}\eqdef \{ x\in


D^{t,1}_{\infty}: v(t)\le x(t)\le z(t),\ t \in [0,1]\}


$$


--- порядковый интервал в пространстве ${ D}^{t,1}_{\infty}$. Пусть


$\Theta: { C}\seta { B}_{\infty}^{t,1}$ --- линейный


изотонный оператор,


$\ov{v}\equiv \Theta v$, $\ov{z}\equiv \Theta z$,


$[\ov{v},\ov{z}]_{B^{t,1}_{\infty}}$


--- порядковый интервал в пространстве


${ B}^{t,1}_{\infty}$.





\begin{defn}[{\series{m}\shape{n}\selectfont \cite{2}, с.\,218}]


Будем говорить, что функция


$f(\cdot,\cdot)$ удовлетворяет


одному из условий


$L^i[\ov{v},\ov{z}]$, $i=1,2,$ если для любой


измеримой функции


$u\in [\ov v,\ov z]_{B^{t,1}_{\infty}}$


при почти всех $t \in [0,1]$ возможно разложение


$f(t,u(t))=q_i(t)u(t)+M_i(t,u(t))$.


Здесь $q_i\in{ L}_\infty$, $i=1,2$,


${\cal M}_i:[\ov{v},\ov{z}]_{B^{t,1}_{\infty}}\seta


{ B}^{t,1}_{\infty}$ ---


оператор Немыцкого, определенный равенством


$({\cal M}_iu)(t)\eqdef M_i(t,u(t))$, причем


оператор ${\cal M}_1$ изотонен,


а оператор ${\cal M}_2$ антитонен.


\end{defn}





Рассмотрим квазилинейную задачу


\begin{equation}


({\cal L}_0 x)(t) \eqdef


\Ddot x(t)+ \frac{\Dot x(t)}{t} - \frac{x(t)}{t^2}=


f(t, (\Theta x)(t)),\quad t\in [0,1],\quad


x(0)=0,\quad x(1) =\alpha,


\label{3}


\end{equation}


где функция


$f(\cdot,\cdot)$ удовлетворяет условиям Каратеодори.





В \cite{4}, \cite{5} встречаются краевые задачи,


которые могут быть представлены в виде (\ref{3}) при


$\Theta x\equiv x.$ В статье \cite{6} указаны условия существования


и единственности решения такой задачи. Ниже сформулируем условия


существования и единственности решения задачи (\ref{3}), которые


непосредственно не следуют из \cite{6}. Задачу (\ref{3})


исследуем в пространстве ${ D}_{\infty}^{t,1}$.





Аналогично теореме 1 из \cite{7} можно показать, что имеет место





\begin{thmnonum}


Пусть существует пара функций


$v, z \in D_{\infty}^{t,1}$, $v(t) < z(t)$,


$t\in (0,1)$, удовлетворяющая неравенствам


$$


{\cal L}_0 z\le f(\cdot, \Theta z),


\ \; {\cal L}_0 v\ge f(\cdot, \Theta v),


\ \; v(1)\le \alpha \le z(1).


$$


Предположим, что функция $f(\cdot,\cdot)$ удовлетворяет условию


$L^2[\ov v, \ov z]$


с коэффициентом


$q_2\in L_\infty$, $q_2(\cdot)\le 0$.


Тогда задача $(\ref{3})$ имеет по крайней мере одно


решение $x\in [v,z]_{D_{\infty}^{t,1}}$.





Если, кроме того, выполнено условие $L^1[\ov v, \ov z]$


с коэффициентом $q_1\in { L}_\infty$,


причем оператор Грина вспомогательной задачи


$$


{\cal L}_1 x\eqdef{\cal L}_0 x - q_1 \Theta x=


\phi,\quad x(1) = 0,


$$


антитонен, то решение задачи $(\ref{3})$ единственно.


\end{thmnonum}





\begin{thebibliography}{9}





\bibitem{1}


Шиндяпин А.И. {\em О краевой задаче для одного сингулярного


уравнения} // Дифференц. уравнения. -- 1984. -- Т.\,20. -- \No\,3. --


С.\,450--455.





\bibitem{2}


Азбелев Н.В., Максимов В.П., Рахматуллина


Л.Ф. {\em Введение в теорию функционально-дифференциальных


уравнений.} -- М.: Наука, 1991. -- 280~с.





\bibitem{3}


Rhoades B.~E. {\em Norm and


spectral properties of some weighted mean operators} //


Math. -- 1984. -- V.\,26. -- \No\,2. --


P.\,143--152.


\bibitem{4}


Срубщик Л.С. {\em Нежесткость сферической оболочки} // ПММ.


-- 1967. -- Т.\,31. -- Вып.\,4. -- С.\,723--729.


\bibitem{5}


Срубщик Л.С. {\em Докритическое равновесие


тонкой пологой оболочки вращения и его устойчивость} // ПММ.


-- 1980. -- Т.\,44. -- Вып.\,2. -- С.\,327--337.


\bibitem{6}


Гризанс Г.П. {\em Об одной краевой задаче для


уравнения с несуммируемой особенностью} // Латв. матем. ежегодник. --


1985. -- Вып.\,29. -- С.\,22--35.


\bibitem{7}


Алвеш М.Ж. О разрешимости двухточечной краевой задачи для


сингулярного нелинейного функционально-дифференциального уравнения //


Изв. вузов. Математика. -- 1999. -- \No\,2. -- С.\,12--19.





\end{thebibliography}





\vskip 2cm





\post{\it Пермский государственный университет}{\it Поступила}


\post{}{24.05.1999}


\label{end}


\end{document}











